%% Annexe
\chapter{Test d'une annexe}

\section{Première Section de l'Annexe}

\subsection{Figures en annexe}

\begin{figure}
    \centering
        \includegraphics[width=0.65\textwidth]{figures/logoESPA.png}
	 \\ \parbox{0.75\textwidth}{\caption{Figure en Annexe.}\label{fig:testAp}}
\end{figure}

Les figures en annexe se déclarent de la même manière que dans le reste du document, et leur numération est automatiquement adaptée (exemple, Figure \ref{fig:testAp}).

\subsubsection{Tables en annexe}

\begin{table}
		\parbox{0.65\textwidth}{\caption{Table en Annexe.}\label{tab:testAp}}

		\begin{tabular}{|c|c|c|c|c|c|c|c|}
		\hline
			{\bf titre} & {\bf titre} & {\bf titre} & {\bf titre} & {\bf titre} & {\bf titre} & {\bf titre} & {\bf titre} \\
	  \hline
			blá & blá & blá & blá & blá & blá & blá & blá \\
	  \hline
			blá & blá & blá & blá & blá & blá & blá & blá \\
	  \hline
			blá & blá & blá & blá & blá & blá & blá & blá \\
	  \hline
			blá & blá & blá & blá & blá & blá & blá & blá \\
	  \hline
			blá & blá & blá & blá & blá & blá & blá & blá \\
	  \hline
			blá & blá & blá & blá & blá & blá & blá & blá \\
	  \hline
		\end{tabular}
\end{table}

Même chose pour les tableaux (exemple, Tableau \ref{tab:testAp}).

\begin{spacing}{0.8}
	\begin{figure*}[!h]
		\begin{center}
			\subfloat[Logo de l'UNA ]{
				\label{comp01-2}
					\includegraphics[width=0.4\textwidth]{figures/logoESPA} %
			}\hspace{1.5cm}%\hfill
			\subfloat[Logo de l'ESPA]{
				\label{comp01-1}
					\includegraphics[width=0.4\textwidth]{figures/logoUNA} %
			}\hfill
			\caption{Deux figures cote a cote}%
			\label{comp01}
		\end{center}
	\end{figure*}
\end{spacing}
